\chapter{Lógica, conjuntos y aplicaciones}\label{ch:b1:logic}

Hasta ahora, las demostraciones se han llevado a cabo con una idea
informal pero honesta de lo que significa «demostrar». Este primer
capítulo de matemáticas universitarias explicita las reglas del juego:
qué es un \omterm{def:b1:logic:statement}{enunciado} matemático, cómo los conectivos y los cuantificadores
combinan \omterm{def:b1:logic:statement}{enunciados}, qué pasos son lícitos en una demostración --- y
construye después, sobre esa base, los dos lenguajes universales de las
matemáticas: los \omterm{def:b1:logic:sets}{conjuntos} y las aplicaciones.

\section{Enunciados y conectivos}

\begin{definition}[Enunciado, conectivos]\label{def:b1:logic:statement}
Un \emph{enunciado}\index{enunciado} (o \emph{proposición}) es una frase
que es verdadera (V) o falsa (F) --- exactamente una de las dos. A
partir de dos enunciados $P$ y $Q$ se forman:
\begin{itemize}
  \item la \emph{negación} $\lnot P$ («no $P$»), verdadera exactamente
        cuando $P$ es falsa;
  \item la \emph{conjunción} $P \land Q$ («$P$ y $Q$»), verdadera
        exactamente cuando ambas lo son;
  \item la \emph{disyunción} $P \lor Q$ («$P$ o $Q$»), verdadera
        exactamente cuando al menos una lo es (este «o» es inclusivo);
  \item la \emph{implicación} $P \implies Q$, falsa exactamente cuando
        $P$ es verdadera y $Q$ falsa;
  \item la \emph{equivalencia} $P \iff Q$, verdadera exactamente cuando
        $P$ y $Q$ tienen el mismo valor de verdad.
\end{itemize}
\end{definition}

\begin{remark}
La tabla de verdad de $P \implies Q$ merece una pausa: cuando $P$ es
falsa, $P \implies Q$ es \emph{verdadera}, sea cual sea $Q$. «Si
$2 < 1$, entonces $0 = 5$» es una implicación verdadera. Una implicación
no afirma nada sobre lo que ocurre cuando su hipótesis falla.
\end{remark}

\begin{proposition}[Reglas de cálculo con enunciados]\label{prop:b1:logic:rules}
Para todos los \omterm{def:b1:logic:statement}{enunciados} $P$, $Q$, $R$:
\begin{enumerate}
  \item $\lnot(\lnot P) \iff P$;
  \item leyes de De Morgan\index{leyes de De Morgan}:
        $\lnot(P \land Q) \iff (\lnot P) \lor (\lnot Q)$ y
        $\lnot(P \lor Q) \iff (\lnot P) \land (\lnot Q)$;
  \item $(P \implies Q) \iff \bigl((\lnot P) \lor Q\bigr)$, de donde
        $\lnot(P \implies Q) \iff P \land (\lnot Q)$;
  \item contraposición\index{contraposición}:
        $(P \implies Q) \iff \bigl((\lnot Q) \implies (\lnot P)\bigr)$;
  \item $(P \iff Q) \iff \bigl((P \implies Q) \land (Q \implies
        P)\bigr)$;
  \item distributividad: $P \land (Q \lor R) \iff (P \land Q) \lor
        (P \land R)$ and $P \lor (Q \land R) \iff (P \lor Q) \land
        (P \lor R)$.
\end{enumerate}
\end{proposition}

\begin{proof}
Cada equivalencia se comprueba comparando tablas de verdad: dos
\omterm{def:b1:logic:statement}{enunciados} compuestos construidos a partir de $P$, $Q$, $R$ son
equivalentes exactamente cuando toman el mismo valor de verdad en cada
uno de los (cuatro u ocho) casos. Escribamos una tabla completa, la de
la primera ley de De Morgan:
\begin{center}
\begin{tabular}{cc|c|c|cc|c}
$P$ & $Q$ & $P \land Q$ & $\lnot(P \land Q)$ & $\lnot P$ & $\lnot Q$ &
$(\lnot P) \lor (\lnot Q)$ \\
\hline
V & V & V & F & F & F & F \\
V & F & F & V & F & V & V \\
F & V & F & V & V & F & V \\
F & F & F & V & V & V & V \\
\end{tabular}
\end{center}
Las columnas $4$ y $7$ coinciden, lo que demuestra la ley. Para la
contraposición resulta más rápido un atajo verbal: $P \implies Q$ es
falsa exactamente en el caso ($P$ verdadera, $Q$ falsa), y
$(\lnot Q) \implies (\lnot P)$ es falsa exactamente en el caso
($\lnot Q$ verdadera, $\lnot P$ falsa), es decir ($Q$ falsa, $P$
verdadera) --- el mismo y único caso, de modo que las dos implicaciones
tienen tablas idénticas. Las demás reglas se comprueban igual;
obsérvese que (3) reduce toda implicación a una disyunción, con lo que
(2) produce mecánicamente la regla de negación $\lnot(P \implies Q)
\iff P \land (\lnot Q)$: para contradecir una implicación hay que
exhibir un caso en el que la hipótesis se cumpla y la conclusión falle.
\end{proof}

\section{Cuantificadores}

\begin{definition}[Cuantificadores]\label{def:b1:logic:quantifiers}
Sea $P(x)$ una propiedad de un elemento $x$ de un \omterm{def:b1:logic:sets}{conjunto} $E$.
\begin{itemize}
  \item $\forall x \in E,\ P(x)$ («para todo $x$ de $E$, $P(x)$») es
        verdadera cuando todo elemento de $E$ satisface $P$;
  \item $\exists x \in E,\ P(x)$ («existe $x$ en $E$ tal que $P(x)$»)
        es verdadera cuando al menos un elemento de $E$ satisface $P$.
\end{itemize}
Se escribe $\exists!$ para «existe un único».
\end{definition}

\begin{proposition}[Negación de los cuantificadores]\label{prop:b1:logic:negquant}
\[
\lnot\bigl(\forall x \in E,\ P(x)\bigr) \iff
\exists x \in E,\ \lnot P(x),
\qquad
\lnot\bigl(\exists x \in E,\ P(x)\bigr) \iff
\forall x \in E,\ \lnot P(x).
\]
\end{proposition}

\begin{proof}
Razonemos la primera equivalencia en los dos sentidos; la segunda es
simétrica. Si $\forall x \in E,\ P(x)$ es falsa, no todo elemento
satisface $P$: el \omterm{def:b1:logic:sets}{conjunto} $A = \{x \in E : \lnot P(x)\}$ no puede ser
vacío, y cualquiera de sus elementos atestigua $\exists x \in
E,\ \lnot P(x)$. Recíprocamente, si algún $x_0 \in E$ cumple $\lnot
P(x_0)$, entonces $x_0$ es un contraejemplo y el \omterm{def:b1:logic:statement}{enunciado} universal
falla. Para la segunda regla: «ningún $x$ satisface $P$» significa que
el \omterm{def:b1:logic:sets}{conjunto} $\{x : P(x)\}$ es vacío, es decir, que todo $x$ está en su
complementario $A$. Aplicadas en cascada a un prefijo de cuantificadores
anidados, las dos reglas dan el procedimiento mecánico del~\cref{ex:b1:logic:limit}:
la negación recorre la fórmula de izquierda a derecha, cambiando cada
$\forall$ por $\exists$ y cada $\exists$ por $\forall$, y niega
finalmente el predicado más interno.
\end{proof}

\begin{example}[Negar frases matemáticas de todos los días]\label{ex:b1:logic:negpractice}
Sea $f \colon \R \to \R$. La frase «$f$ es creciente» se escribe
\[
\forall x \in \R,\ \forall y \in \R,\quad
x \leq y \implies f(x) \leq f(y) ,
\]
y su negación, por la~\cref{prop:b1:logic:negquant} junto con la regla
$\lnot(P \implies Q) \iff P \land \lnot Q$:
\[
\exists x \in \R,\ \exists y \in \R,\quad
x \leq y \ \text{ y }\ f(x) > f(y) :
\]
basta con un par que lo atestigüe. Del mismo modo, «$f$ está acotada»
es $\exists M \in \R,\ \forall x \in \R,\ \abs{f(x)} \leq M$, con
negación
\[
\forall M \in \R,\ \exists x \in \R,\quad \abs{f(x)} > M :
\]
\emph{sea cual sea} la cota propuesta, algún punto la supera. La idea
clave: una negación correcta nunca contiene un «no» aplicado a un bloque
cuantificado --- es un nuevo \omterm{def:b1:logic:statement}{enunciado} positivo en el que los papeles se
intercambian: ahora uno produce los testigos que antes recibía.
\end{example}

\begin{example}[Orden de los cuantificadores]\label{ex:b1:logic:order}
El \omterm{def:b1:logic:order}{orden} de dos cuantificadores distintos importa:
\[
\forall x \in \R,\ \exists y \in \R,\ y > x
\quad\text{es verdadera (tómese } y = x+1\text{),}
\]
\[
\exists y \in \R,\ \forall x \in \R,\ y > x
\quad\text{es falsa (ningún real supera a todos los reales).}
\]
En el primer \omterm{def:b1:logic:statement}{enunciado}, $y$ puede depender de $x$; en el segundo, un
único $y$ debe servir para todo $x$. Dos cuantificadores iguales, en
cambio, siempre conmutan.
\end{example}

\begin{example}[Leer una definición con tres cuantificadores]\label{ex:b1:logic:limit}
La frase «la sucesión $(u_n)$ converge a $\ell$» se escribirá en
el~\cref{ch:b1:seq} como
\[
\forall \varepsilon > 0,\ \exists N \in \N,\ \forall n \geq N,\quad
\abs{u_n - \ell} \leq \varepsilon .
\]
Su negación, aplicando tres veces la~\cref{prop:b1:logic:negquant}, es
\[
\exists \varepsilon > 0,\ \forall N \in \N,\ \exists n \geq N,\quad
\abs{u_n - \ell} > \varepsilon .
\]
Saber negar mecánicamente frases de este tipo, sin pensar en lo que
significan, es una destreza real: separa el trabajo lógico del trabajo
matemático.
\end{example}

\section{Técnicas de demostración}

\begin{method}[Los esquemas de demostración habituales]\label{met:b1:logic:proofs}
Para demostrar\dots
\begin{enumerate}
  \item \emph{una implicación $P \implies Q$ directamente}: se supone
        $P$ y se deduce $Q$;
  \item \emph{por contraposición}: se supone $\lnot Q$ y se deduce
        $\lnot P$ --- válido por la~\cref{prop:b1:logic:rules}~(4);
  \item \emph{por reducción al absurdo}: se supone que el \omterm{def:b1:logic:statement}{enunciado} es
        falso y se deriva una contradicción;
  \item \emph{una equivalencia}: se demuestran las dos implicaciones por
        separado (o se encadenan equivalencias conocidas);
  \item \emph{un \omterm{def:b1:logic:statement}{enunciado} «para todo»}: se toma un $x$
        \emph{arbitrario} de $E$ («sea $x \in E$») y se demuestra $P(x)$;
  \item \emph{un \omterm{def:b1:logic:statement}{enunciado} «existe»}: se exhibe un testigo, o se
        demuestra la existencia de forma indirecta;
  \item \emph{por inducción}: véase el~\cref{thm:b1:logic:induction}.
\end{enumerate}
Al demostrar un \omterm{def:b1:logic:statement}{enunciado} sobre un elemento bien elegido pero
arbitrario, nunca hay que atribuirle propiedades adicionales: «sea
$x \in \R$» seguido de «como $x > 0$\dots» no demuestra nada sobre los
$x$ negativos.
\end{method}

\begin{remark}[Errores frecuentes en las demostraciones]\label{rem:b1:logic:pitfalls}
Cuatro trampas clásicas, que conviene nombrar de una vez.
\begin{enumerate}
  \item \emph{El recíproco en lugar del contrarrecíproco.}
        $Q \implies P$ \emph{no} equivale a $P \implies Q$; solo lo hace
        $\lnot Q \implies \lnot P$. De «si llueve, la calle se moja» no
        se puede concluir que haya llovido porque la calle esté mojada.
  \item \emph{Demostrar una equivalencia con una sola implicación.} Una
        afirmación «si y solo si» son dos teoremas; hay que anunciar qué
        sentido se demuestra y demostrar los dos. Una cadena de $\iff$
        solo es lícita si \emph{todos} sus eslabones son realmente
        reversibles --- elevar al cuadrado una ecuación, por ejemplo, no
        lo es.
  \item \emph{Demostraciones hacia atrás.} Partir de la conclusión
        deseada y deducir un \omterm{def:b1:logic:statement}{enunciado} verdadero no demuestra nada (de
        $-1 = 1$ se deduce, elevando al cuadrado, el verdadero
        $1 = 1$). Un cálculo puede \emph{descubrirse} hacia atrás, pero
        debe \emph{escribirse} hacia adelante, o con equivalencias
        explícitas.
  \item \emph{Testigo fijo frente a elemento arbitrario.} Para demostrar
        $\exists x,\ P(x)$ basta exhibir un $x$ hábilmente elegido; para
        demostrar $\forall x,\ P(x)$, el $x$ elegido debe seguir siendo
        arbitrario. Confundir ambas cosas --- comprobar una afirmación
        universal en un ejemplo --- es el error más frecuente entre
        quienes empiezan.
\end{enumerate}
\end{remark}

\begin{example}[Contraposición y absurdo en acción]\label{ex:b1:logic:sqrt2}
\emph{Para $n \in \N$: si $n^2$ es par, entonces $n$ es par.} Por
contraposición: si $n$ es impar, $n = 2k+1$, entonces
$n^2 = 4k^2 + 4k + 1$ es impar.

\emph{$\sqrt 2$ es irracional.} Por reducción al absurdo: supongamos que
$\sqrt 2 = p/q$ con $p, q \in \N^*$ y la fracción irreducible. Entonces
$p^2 = 2q^2$ es par, luego $p$ es par (punto anterior), $p = 2r$;
entonces $q^2 = 2r^2$ es par, luego $q$ es par --- lo que contradice la
irreducibilidad.
\end{example}

\begin{theorem}[Inducción]\label{thm:b1:logic:induction}
Sea $P(n)$ una propiedad del entero $n$. Si
\begin{enumerate}
  \item $P(0)$ es verdadera, y
  \item para todo $n \in \N$, $P(n) \implies P(n+1)$,
\end{enumerate}
entonces $P(n)$ es verdadera para todo $n \in \N$.

\emph{Inducción fuerte:} la conclusión no cambia si se sustituye (2)
por: para todo $n$, $\bigl(P(0) \land \dots \land P(n)\bigr) \implies
P(n+1)$.
\end{theorem}

\begin{proof}
Se trata de una propiedad del propio $\N$, equivalente a esta:
\emph{todo subconjunto no vacío de $\N$ tiene un elemento mínimo} (que
damos por conocida). En efecto, supongamos (1) y (2) y sea
$A = \{n \in \N : P(n) \text{ falsa}\}$. Si $A \neq \emptyset$, tiene un
elemento mínimo $m$; $m \neq 0$ por (1); entonces $m - 1 \notin A$,
luego $P(m-1)$ se cumple, y (2) da $P(m)$ --- contradicción. Por tanto
$A = \emptyset$. Para la inducción fuerte se aplica el mismo
razonamiento: $P(0), \dots, P(m-1)$ se cumplen todas, pues $m$ es el
mínimo de $A$.
\end{proof}

\begin{example}[Demostrar existencia y unicidad]\label{ex:b1:logic:uniqueexistence}
Un \omterm{def:b1:logic:statement}{enunciado} $\exists!\,x,\ P(x)$ son \emph{dos} \omterm{def:b1:logic:statement}{enunciados}, que se
demuestran por separado: la existencia (exhibir o construir algún $x_0$
con $P(x_0)$) y la unicidad (suponer $P(x)$ y $P(x')$ y deducir
$x = x'$). Muestra: \emph{existe un único real $x$ tal que $x^3 + x =
2$.} Existencia: $x_0 = 1$ sirve, pues $1 + 1 = 2$. Unicidad: si
$x^3 + x = x'^3 + x'$, entonces
\[
0 = (x^3 - x'^3) + (x - x')
= (x - x')\,\bigl(x^2 + xx' + x'^2 + 1\bigr),
\]
y el segundo factor es positivo (vale $\bigl(x +
\tfrac{x'}2\bigr)^2 + \tfrac34 x'^2 + 1 \geq 1$), luego $x = x'$.
Obsérvese el reparto del trabajo: la existencia se apoyó en una
conjetura afortunada; la unicidad, en un cálculo algebraico válido para
soluciones \emph{arbitrarias} --- ninguno de los dos argumentos hace el
trabajo del otro, y olvidar la segunda mitad es una tentación constante
en cuanto se ha encontrado una solución.
\end{example}

\begin{example}\label{ex:b1:logic:induction}
Para todo $n \in \N^*$: $\;\sum_{k=1}^n k = \frac{n(n+1)}{2}$. Caso base
$n = 1$: los dos miembros valen $1$. Paso: suponiendo la fórmula para $n$,
\[
\sum_{k=1}^{n+1} k = \frac{n(n+1)}{2} + (n+1)
= (n+1)\Bigl(\frac n2 + 1\Bigr) = \frac{(n+1)(n+2)}{2}. \qedhere
\]
\end{example}

\begin{example}[La inducción fuerte en acción]\label{ex:b1:logic:strongind}
\emph{Todo entero $n \geq 2$ es un producto de números primos} (siendo
un primo un entero $\geq 2$ cuyos únicos divisores $\geq 1$ son $1$ y él
mismo; los primos se estudian por sí mismos en el~\cref{ch:b1:arith}).
La inducción ordinaria es impotente aquí: saber que $95 = 5 \times 19$
se factoriza no dice nada sobre $96$. La inducción fuerte encaja
exactamente. Caso base: $2$ es primo, luego es un producto (de un solo
factor) de primos. Paso: sea $n \geq 2$ y supongamos que todo entero $m$
con $2 \leq m \leq n$ es un producto de primos. Si $n + 1$ es primo, ya
está. En caso contrario $n + 1 = ab$ con $2 \leq a, b \leq n$; por la
hipótesis fuerte, $a$ y $b$ son productos de primos, y por tanto también
lo es $n + 1$. La idea clave: la inducción fuerte es la herramienta
adecuada siempre que la «razón» de $P(n+1)$ resida en un rango anterior
imprevisible, y no en el rango $n$.
\end{example}

\section{Conjuntos}

\begin{definition}[Operaciones con conjuntos]\label{def:b1:logic:sets}
Tomamos como primitivas la noción de \emph{conjunto}\index{conjunto} y
la relación de pertenencia $x \in E$. Para conjuntos $A, B$ contenidos
en un conjunto ambiente $E$:
\begin{itemize}
  \item \emph{inclusión}: $A \subseteq B$ cuando $\forall x,\ x \in A
        \implies x \in B$; igualdad $A = B$ cuando $A \subseteq B$ y
        $B \subseteq A$;
  \item \emph{unión} $A \cup B$, \emph{intersección} $A \cap B$,
        \emph{diferencia} $A \setminus B = \{x \in A : x \notin B\}$,
        \emph{complementario} $\overline{A} = E \setminus A$;
  \item el \emph{conjunto vacío} $\emptyset$, contenido en todo conjunto;
  \item el \emph{conjunto de las partes}\index{conjunto de las partes}
        $\mathcal{P}(E)$: el conjunto de todos los subconjuntos de $E$;
  \item el \emph{producto} $E \times F$: el conjunto de los pares
        ordenados $(x, y)$ con $x \in E$, $y \in F$.
\end{itemize}
\end{definition}

\begin{example}[Familiarizarse con el conjunto de las partes]\label{ex:b1:logic:powerset}
Para $E = \{a, b\}$:
\[
\mathcal P(E) = \bigl\{\, \emptyset,\ \{a\},\ \{b\},\ \{a, b\}
\,\bigr\},
\]
cuatro elementos --- y obsérvese la disciplina de tipos: $a \in E$, pero
$\{a\} \in \mathcal P(E)$; los \omterm{def:b1:logic:statement}{enunciados} $a \in \mathcal P(E)$
y $\{a\} \subseteq \mathcal P(E)$ son ambos falsos tal como están
escritos (el segundo exigiría que $a$ fuese un \emph{subconjunto} de
$E$). Iterando desde la nada: $\mathcal P(\emptyset) = \{\emptyset\}$
tiene un elemento, $\mathcal P(\mathcal P(\emptyset)) = \{\emptyset,
\{\emptyset\}\}$ tiene dos, el siguiente tiene cuatro --- los \omterm{def:b1:logic:sets}{conjuntos}
de \omterm{def:b1:logic:sets}{conjuntos} son \omterm{def:b1:logic:sets}{conjuntos} ordinarios, y el~\cref{ch:b1:counting}
confirmará la duplicación: $\abs{\mathcal P(E)} = 2^{\abs E}$. Mantener
claros los niveles ($x$, $\{x\}$, $\{\{x\}\}$) es la mitad del trabajo
en ejercicios como los~\cref{exo:b1:logic:11,exo:b1:logic:12}.
\end{example}

\begin{proposition}[Álgebra de conjuntos]\label{prop:b1:logic:setalgebra}
Para subconjuntos $A, B, C$ de $E$:
\begin{enumerate}
  \item $A \cap (B \cup C) = (A \cap B) \cup (A \cap C)$ y
        $A \cup (B \cap C) = (A \cup B) \cap (A \cup C)$;
  \item De Morgan: $\overline{A \cup B} = \overline{A} \cap
        \overline{B}$ y $\overline{A \cap B} = \overline{A} \cup
        \overline{B}$;
  \item $A \subseteq B \iff \overline{B} \subseteq \overline{A}$.
\end{enumerate}
\end{proposition}

\begin{proof}
Cada identidad traduce una regla de la~\cref{prop:b1:logic:rules}
mediante el diccionario ($\in A$ o no) $\leftrightarrow$ (\omterm{def:b1:logic:statement}{enunciado}
verdadero o falso): por ejemplo, $x \in \overline{A \cup B} \iff
\lnot(x \in A \lor x \in B) \iff (x \notin A) \land (x \notin B) \iff
x \in \overline{A} \cap \overline{B}$. El punto (3) es la contraposición.
Como segunda muestra, la primera ley distributiva completa:
\[
x \in A \cap (B \cup C)
\iff (x \in A) \land \bigl(x \in B \lor x \in C\bigr)
\iff \bigl(x \in A \land x \in B\bigr) \lor
\bigl(x \in A \land x \in C\bigr),
\]
por la distributividad de la~\cref{prop:b1:logic:rules}~(6), y el último
\omterm{def:b1:logic:statement}{enunciado} se lee $x \in (A \cap B) \cup (A \cap C)$. Toda identidad
conjuntista de este tipo se demuestra con esta única traducción
mecánica --- razón por la cual ninguna de ellas hay que memorizarla.
\end{proof}

\begin{method}[Demostrar igualdades de conjuntos]\label{met:b1:logic:doubleinc}
Para demostrar $A = B$ se prueban las dos inclusiones: sea $x \in A$, se
comprueba que $x \in B$; después, sea $x \in B$, se comprueba que
$x \in A$. Otra vía es encadenar equivalencias $x \in A \iff \dots \iff
x \in B$, cuando cada paso sea realmente una equivalencia.
\end{method}

\begin{omfigure}
\begin{tikzpicture}[scale=0.9]
  \draw[thick] (-2.6,-1.9) rectangle (2.6,1.9);
  \node[below right] at (-2.6,1.9) {$E$};
  \begin{scope}
    \clip (-2.6,-1.9) rectangle (2.6,1.9);
    \fill[omDef!20] (-2.6,-1.9) rectangle (2.6,1.9);
    \fill[white] (-0.7,0) circle (1.1);
    \fill[white] (0.7,0) circle (1.1);
  \end{scope}
  \draw[omDef, thick] (-0.7,0) circle (1.1);
  \draw[omDef, thick] (0.7,0) circle (1.1);
  \node[omDef] at (-1.35,0.5) {$A$};
  \node[omDef] at (1.35,0.5) {$B$};
  \begin{scope}[xshift=6.6cm]
    \draw[thick] (-2.6,-1.9) rectangle (2.6,1.9);
    \node[below right] at (-2.6,1.9) {$E$};
    \begin{scope}
      \clip (-2.6,-1.9) rectangle (2.6,1.9);
      \fill[omThm!20] (-2.6,-1.9) rectangle (2.6,1.9);
      \begin{scope}
        \clip (-0.7,0) circle (1.1);
        \fill[white] (0.7,0) circle (1.1);
      \end{scope}
    \end{scope}
    \draw[omThm, thick] (-0.7,0) circle (1.1);
    \draw[omThm, thick] (0.7,0) circle (1.1);
    \node[omThm] at (-1.35,0.5) {$A$};
    \node[omThm] at (1.35,0.5) {$B$};
  \end{scope}
\end{tikzpicture}

{\small Las leyes de De Morgan en imágenes: la región sombreada de la
izquierda es $\overline{A \cup B} = \overline A \cap \overline B$ (todo
lo que queda fuera de los dos discos); a la derecha,
$\overline{A \cap B} = \overline A \cup \overline B$ (todo salvo la
lente central). Un dibujo no es una demostración, pero hace imposible
recordar mal la demostración por elementos de
la~\cref{prop:b1:logic:setalgebra}.}
\end{omfigure}

\section{Aplicaciones}

\begin{definition}[Aplicación, imagen, imagen recíproca]\label{def:b1:logic:map}
Una \emph{aplicación}\index{aplicación} (o \emph{función})
$f \colon E \to F$ asigna a cada elemento $x$ del \omterm{def:b1:logic:sets}{conjunto} $E$ (el
\emph{dominio}) exactamente un elemento $f(x)$ del \omterm{def:b1:logic:sets}{conjunto} $F$ (el
\emph{codominio}). Para $A \subseteq E$ y $B \subseteq F$:
\[
f(A) = \{f(x) : x \in A\} \subseteq F,
\qquad
f^{-1}(B) = \{x \in E : f(x) \in B\} \subseteq E
\]
son la \emph{imagen directa} de $A$ y la \emph{imagen recíproca}
\index{imagen recíproca} de $B$. La \emph{composición} de
$f \colon E \to F$ y $g \colon F \to G$ es $g \circ f \colon E \to G$,
$x \mapsto g(f(x))$.
\end{definition}

\begin{remark}
La notación $f^{-1}(B)$ \emph{no} presupone que exista una \omterm{def:b1:logic:map}{aplicación}
inversa: $f^{-1}(B)$ está definida para toda $f$. Las imágenes
recíprocas se comportan mejor que las directas: $f^{-1}$ conserva
uniones, intersecciones y complementarios, mientras que
$f(A \cap A') \subseteq f(A) \cap f(A')$ puede ser estricta
(\cref{exo:b1:logic:8}).
\end{remark}

\begin{example}[Cálculo de imágenes e imágenes recíprocas]\label{ex:b1:logic:images}
Sea $f \colon \R \to \R$, $x \mapsto x^2$. Entonces:
\[
f\bigl(\intcc{-1}{2}\bigr) = \intcc04, \qquad
f^{-1}\bigl(\intcc14\bigr) = \intcc{-2}{-1} \cup \intcc12, \qquad
f^{-1}(\{-1\}) = \emptyset .
\]
Para la primera: todo $x \in \intcc{-1}2$ cumple $x^2 \in \intcc04$, y
todo $y \in \intcc04$ se alcanza como $y = (\sqrt y)^2$ con $\sqrt y
\in \intcc02 \subseteq \intcc{-1}2$ --- obsérvese que la imagen
\emph{no} es $\intcc14 = \{(-1)^2, 2^2\}$: las imágenes de intervalos no
se calculan solo con los extremos. Para la segunda: $1 \leq x^2 \leq 4
\iff 1 \leq \abs x \leq 2$, que se parte en dos trozos. La tercera
ilustra que una \omterm{def:b1:logic:map}{imagen recíproca} puede ser vacía --- $f^{-1}(B)$
siempre tiene sentido, por pequeña que sea la intersección de $B$ con la
imagen. Por último, obsérvese en este ejemplo el fenómeno de inclusión
estricta de la observación anterior: con $A = \intcc{-1}0$ y
$A' = \intcc01$ se tiene $f(A \cap A') = f(\{0\}) = \{0\}$, mientras que
$f(A) \cap f(A') = \intcc01$.
\end{example}

\begin{definition}[Inyectiva, sobreyectiva, biyectiva]\label{def:b1:logic:inj}
Una \omterm{def:b1:logic:map}{aplicación} $f \colon E \to F$ es:
\begin{itemize}
  \item \emph{inyectiva}\index{inyectiva} cuando elementos distintos
        tienen imágenes distintas: $\forall x, x' \in E,\ f(x) = f(x')
        \implies x = x'$;
  \item \emph{sobreyectiva}\index{sobreyectiva} cuando todo elemento de
        $F$ se alcanza: $\forall y \in F,\ \exists x \in E,\ f(x) = y$;
  \item \emph{biyectiva}\index{biyectiva} cuando es ambas cosas, es
        decir, cuando todo $y \in F$ tiene exactamente una
        \omterm{def:b1:logic:map}{imagen recíproca}.
\end{itemize}
\end{definition}

\begin{theorem}[Aplicación inversa]\label{thm:b1:logic:inverse}
Una \omterm{def:b1:logic:map}{aplicación} $f \colon E \to F$ es \omterm{def:b1:logic:inj}{biyectiva} si y solo si existe
una \omterm{def:b1:logic:map}{aplicación} $g \colon F \to E$ tal que $g \circ f = \mathrm{id}_E$
y $f \circ g = \mathrm{id}_F$. En tal caso $g$ es única; se escribe
$f^{-1}$ y se llama \emph{inversa} de $f$, y $f^{-1}$ es a su vez
\omterm{def:b1:logic:inj}{biyectiva}, con $(f^{-1})^{-1} = f$.
\end{theorem}

\begin{proof}
($\Rightarrow$) Si $f$ es \omterm{def:b1:logic:inj}{biyectiva}, todo $y \in F$ tiene una única
\omterm{def:b1:logic:map}{imagen recíproca}; se define $g(y)$ como esa \omterm{def:b1:logic:map}{imagen recíproca}.
Entonces $f(g(y)) = y$ por construcción, y $g(f(x)) = x$ porque $x$ es
\emph{la} \omterm{def:b1:logic:map}{imagen recíproca} de $f(x)$.

($\Leftarrow$) Supongamos que existe tal $g$. Si $f(x) = f(x')$,
aplicando $g$ se obtiene $x = x'$: $f$ es \omterm{def:b1:logic:inj}{inyectiva}. Para $y \in F$,
$x = g(y)$ cumple $f(x) = y$: $f$ es \omterm{def:b1:logic:inj}{sobreyectiva}.

Unicidad: si $g$ y $h$ sirven las dos, entonces $g = g \circ
\mathrm{id}_F = g \circ (f \circ h) = (g \circ f) \circ h = h$. Por
último, el par de identidades es simétrico en $f$ y $g$, luego
$g = f^{-1}$ es \omterm{def:b1:logic:inj}{biyectiva} con inversa $f$.
\end{proof}

\begin{example}[Cálculo de una inversa en la práctica]\label{ex:b1:logic:inversecomp}
Sea $f \colon \R \to \intoo0{+\infty}$, $f(x) = \eu^{2x+1}$. Para
invertirla se resuelve $y = f(x)$ para un $y > 0$ dado:
\[
y = \eu^{2x+1} \iff \ln y = 2x + 1 \iff x = \frac{\ln y - 1}2 ,
\]
siendo cada paso reversible en los dominios anunciados. El cálculo lo
entrega todo a la vez: para cada $y$ del codominio hay exactamente una
solución $x$, luego $f$ es \omterm{def:b1:logic:inj}{biyectiva}, y
\[
f^{-1} \colon \intoo0{+\infty} \to \R,
\qquad
f^{-1}(y) = \frac{\ln y - 1}2 .
\]
Una comprobación rápida de las dos composiciones ($f^{-1}(f(x)) =
\frac{(2x+1) - 1}2 = x$ y $f(f^{-1}(y)) = \eu^{\ln y} = y$) confirma el
criterio del~\cref{thm:b1:logic:inverse}. La idea clave: «despejar $x$ y
vigilar las equivalencias» es a la vez la demostración de existencia, la
de unicidad y la fórmula --- pero solo funciona si el codominio se
anunció correctamente ($f$ \emph{no} es \omterm{def:b1:logic:inj}{sobreyectiva} sobre $\R$).
\end{example}

\begin{proposition}[La composición y las tres propiedades]\label{prop:b1:logic:comp}
Sean $f \colon E \to F$ y $g \colon F \to G$.
\begin{enumerate}
  \item Si $f$ y $g$ son \omterm{def:b1:logic:inj}{inyectivas} (resp.\ \omterm{def:b1:logic:inj}{sobreyectivas},
        \omterm{def:b1:logic:inj}{biyectivas}), también lo es $g \circ f$; y entonces
        $(g \circ f)^{-1} = f^{-1} \circ g^{-1}$ en el caso
        \omterm{def:b1:logic:inj}{biyectivo}.
  \item Si $g \circ f$ es \omterm{def:b1:logic:inj}{inyectiva}, entonces $f$ es \omterm{def:b1:logic:inj}{inyectiva}. Si
        $g \circ f$ es \omterm{def:b1:logic:inj}{sobreyectiva}, entonces $g$ es \omterm{def:b1:logic:inj}{sobreyectiva}.
\end{enumerate}
\end{proposition}

\begin{proof}
(1) Si $g(f(x)) = g(f(x'))$, la \omterm{def:b1:logic:inj}{inyectividad} de $g$ da $f(x) = f(x')$, y
la de $f$ da entonces $x = x'$. Si $z \in G$, la \omterm{def:b1:logic:inj}{sobreyectividad} de $g$
da un $y$ con $g(y) = z$, y la de $f$ da un $x$ con $f(x) = y$, de modo
que $g(f(x)) = z$. En el caso \omterm{def:b1:logic:inj}{biyectivo} se comprueba directamente que
$f^{-1} \circ g^{-1}$ es una inversa por los dos lados de $g \circ f$, y
la unicidad del~\cref{thm:b1:logic:inverse} concluye.

(2) Si $f(x) = f(x')$, entonces $g(f(x)) = g(f(x'))$, y la \omterm{def:b1:logic:inj}{inyectividad}
de $g \circ f$ da $x = x'$. Si $z \in G$, la \omterm{def:b1:logic:inj}{sobreyectividad} de
$g \circ f$ da un $x$ con $g(f(x)) = z$: entonces $y = f(x)$ cumple
$g(y) = z$.
\end{proof}

\begin{example}[El punto (2) no se puede mejorar]\label{ex:b1:logic:compsharp}
En la~\cref{prop:b1:logic:comp}~(2) no se pueden reforzar las
conclusiones: que $g \circ f$ sea \omterm{def:b1:logic:inj}{biyectiva} \emph{no} obliga a que $f$
sea \omterm{def:b1:logic:inj}{sobreyectiva} ni a que $g$ sea \omterm{def:b1:logic:inj}{inyectiva}. Tómese $E = G = \{1\}$,
$F = \{1, 2\}$, con $f(1) = 1$ y $g(1) = g(2) = 1$: entonces
$g \circ f = \mathrm{id}_E$ es \omterm{def:b1:logic:inj}{biyectiva}, pero $f$ no alcanza el
elemento $2$ y $g$ colapsa los dos elementos. La moraleja es una regla
de contabilidad precisa: la información sobre la composición pasa a la
\omterm{def:b1:logic:map}{aplicación} \emph{interior} para la \omterm{def:b1:logic:inj}{inyectividad} y a la \emph{exterior}
para la \omterm{def:b1:logic:inj}{sobreyectividad}, nunca al revés. (El~\cref{exo:b1:logic:9}
construye el mismo fenómeno con \omterm{def:b1:logic:sets}{conjuntos} infinitos, donde es el motor
de las inversas laterales.)
\end{example}

\begin{example}\label{ex:b1:logic:injexamples}
$f \colon \R \to \R$, $x \mapsto x^2$ no es \omterm{def:b1:logic:inj}{inyectiva} ($f(-1) =
f(1)$) ni \omterm{def:b1:logic:inj}{sobreyectiva} ($-1$ no tiene \omterm{def:b1:logic:map}{imagen recíproca}). Restringiendo
el dominio y el codominio, $f \colon \R_+ \to \R_+$, $x \mapsto x^2$ sí
es \omterm{def:b1:logic:inj}{biyectiva}, con inversa $y \mapsto \sqrt y$. Que una \omterm{def:b1:logic:map}{aplicación}
sea \omterm{def:b1:logic:inj}{inyectiva} o \omterm{def:b1:logic:inj}{sobreyectiva} depende del dominio y del codominio
anunciados, no solo de la fórmula.
\end{example}

\section{Relaciones}

\begin{definition}[Relación de equivalencia]\label{def:b1:logic:equiv}
Una \emph{relación binaria} $\mathcal{R}$ en un \omterm{def:b1:logic:sets}{conjunto} $E$ es una
\emph{relación de equivalencia}\index{relación de equivalencia} cuando
es \emph{reflexiva} ($x \mathbin{\mathcal{R}} x$ para todo $x$),
\emph{simétrica} ($x \mathbin{\mathcal{R}} y \implies y
\mathbin{\mathcal{R}} x$) y \emph{transitiva} ($x
\mathbin{\mathcal{R}} y$ e $y \mathbin{\mathcal{R}} z$ implican $x
\mathbin{\mathcal{R}} z$). La \emph{clase de equivalencia} de $x$ es
$\mathrm{cl}(x) = \{y \in E : x \mathbin{\mathcal{R}} y\}$.
\end{definition}

\begin{example}[Comprobar los tres axiomas]\label{ex:b1:logic:fracpart}
En $\R$, se declara $x \mathbin{\mathcal{R}} y$ cuando $x - y \in \Z$.
\emph{Reflexiva:} $x - x = 0 \in \Z$. \emph{Simétrica:} si $x - y
\in \Z$, entonces $y - x = -(x - y) \in \Z$. \emph{Transitiva:} si $x -
y \in \Z$ e $y - z \in \Z$, entonces $x - z = (x - y) + (y - z) \in
\Z$ (suma de enteros). Así pues, $\mathcal R$ es una \omterm{def:b1:logic:equiv}{relación de
equivalencia}, y $\mathrm{cl}(x) = x + \Z = \{x + k : k \in \Z\}$: cada
clase contiene exactamente un representante en $\intco01$, su
\emph{parte fraccionaria}. En cambio, la relación «$\abs{x - y}
\leq 1$» en $\R$ es reflexiva y simétrica, pero \emph{no} transitiva
($0 \mathbin{\mathcal R} 1$ y $1 \mathbin{\mathcal R} 2$, y sin embargo
$\abs{0 - 2} > 1$): la proximidad no se propaga, y no existe ninguna
\omterm{thm:b1:logic:partition}{partición} en clases --- un contraejemplo útil cuando comprobar los
axiomas empieza a parecer rutinario.
\end{example}

\begin{theorem}[Las clases forman una partición]\label{thm:b1:logic:partition}
Sea $\mathcal{R}$ una \omterm{def:b1:logic:equiv}{relación de equivalencia} en $E$. Entonces las
\omterm{def:b1:logic:equiv}{clases de equivalencia} son no vacías, dos a dos disjuntas o iguales, y
su unión es $E$: forman una \emph{partición}\index{partición} de $E$.
Recíprocamente, toda \omterm{thm:b1:logic:partition}{partición} de $E$ proviene de este modo de
exactamente una \omterm{def:b1:logic:equiv}{relación de equivalencia} («estar en el mismo trozo»).
\end{theorem}

\begin{proof}
$x \in \mathrm{cl}(x)$ por reflexividad, luego las clases son no vacías y
su unión es $E$. Supongamos que $\mathrm{cl}(x) \cap \mathrm{cl}(y) \neq
\emptyset$, digamos que $z$ está en las dos. Entonces
$x \mathbin{\mathcal{R}} z$ e $y \mathbin{\mathcal{R}} z$, luego, por
simetría y transitividad, $x \mathbin{\mathcal{R}} y$. Ahora bien, para
cualquier $t \in \mathrm{cl}(y)$, la transitividad da
$t \in \mathrm{cl}(x)$, y simétricamente: las dos clases son iguales.
Para el recíproco, sea $(E_i)_{i \in I}$ una \omterm{thm:b1:logic:partition}{partición} de $E$ y
definamos $x \mathbin{\mathcal S} y$ como «algún trozo contiene a la vez
a $x$ y a $y$». \emph{Reflexiva:} $x$ está en algún trozo, que
contiene entonces a $x$ dos veces. \emph{Simétrica:} la condición que
la define es simétrica en $x$ e $y$. \emph{Transitiva:} si
$x, y \in E_i$ e $y, z \in E_j$, entonces $y \in E_i \cap E_j$, luego
$E_i = E_j$ (dos trozos distintos son disjuntos) y $x, z$ comparten
trozo. La clase de $x$ para $\mathcal S$ es exactamente el trozo que
contiene a $x$, de modo que las clases son los trozos dados. Por último,
la relación queda determinada por sus clases: dos \omterm{def:b1:logic:equiv}{relaciones de
equivalencia} con las mismas clases relacionan los mismos pares, ya que
cada una relaciona $x$ e $y$ exactamente cuando $y$ pertenece a la clase
de $x$ --- de donde la unicidad afirmada.
\end{proof}

\begin{example}\label{ex:b1:logic:congruence}
En $\Z$, la congruencia módulo $n$ ($x \equiv y \pmod n$ cuando $n$
divide a $x - y$) es una \omterm{def:b1:logic:equiv}{relación de equivalencia}; sus clases son los
$n$ \omterm{def:b1:logic:sets}{conjuntos} de enteros con un resto dado en la división por $n$.
Este ejemplo se convierte en el anillo $\Z/n\Z$ en el~\cref{ch:b1:structures}.
\end{example}

\begin{definition}[Relación de orden]\label{def:b1:logic:order}
Una relación $\preceq$ en $E$ es un \emph{orden}\index{relación de orden}
cuando es reflexiva, \emph{antisimétrica} ($x \preceq y$ e $y \preceq x$
implican $x = y$) y transitiva. El orden es \emph{total} cuando dos
elementos cualesquiera son comparables, y \emph{parcial} en caso
contrario. Un elemento $M \in A \subseteq E$ es un \emph{máximo} de $A$
cuando $a \preceq M$ para todo $a \in A$; el máximo (y el mínimo) es
único cuando existe.
\end{definition}

\begin{example}\label{ex:b1:logic:orders}
$(\R, \leq)$ está totalmente ordenado. $(\mathcal{P}(E), \subseteq)$
está parcialmente ordenado en cuanto $E$ tiene dos elementos: $\{a\}$ y
$\{b\}$ no son comparables. El subconjunto $A = \{\{a\}, \{b\}\}$ de
$\mathcal{P}(\{a,b\})$ no tiene máximo, y sin embargo tiene una cota
superior, $\{a, b\}$: la distinción entre máximos y cotas superiores
reaparece, para $\R$, en el~\cref{ch:b1:reals}.
\end{example}

\begin{example}[Dos órdenes en la cuadrícula $\N^2$]\label{ex:b1:logic:lexorder}
En los pares de naturales, se comparan componente a componente:
$(a, b) \preceq (a', b')$ cuando $a \leq a'$ \emph{y} $b \leq b'$ (el
\emph{\omterm{def:b1:logic:order}{orden} producto}). Es un \omterm{def:b1:logic:order}{orden} --- cada axioma se hereda
coordenada a coordenada --- pero parcial: $(1, 3)$ y $(2, 0)$ son
incomparables. Comparémoslos ahora como en un diccionario:
$(a, b) \preceq_{\mathrm{lex}} (a', b')$ cuando $a < a'$, o bien
$a = a'$ y $b \leq b'$ (el \emph{\omterm{def:b1:logic:order}{orden} lexicográfico}). La
transitividad exige distinguir dos casos, pero se cumple, y ahora dos
pares cualesquiera son comparables: el \omterm{def:b1:logic:order}{orden} es total. Los dos
órdenes ordenan el mismo \omterm{def:b1:logic:sets}{conjunto} de manera distinta ---
$(0, 100) \preceq_{\mathrm{lex}} (1, 0)$ aunque el \omterm{def:b1:logic:order}{orden} producto no
diga nada --- lo que recuerda que un \omterm{def:b1:logic:order}{orden} es una estructura que se
\emph{elige}, no una propiedad del \omterm{def:b1:logic:sets}{conjunto}. La comparación
lexicográfica es además el truco habitual para reducir varios criterios
de ordenación a uno solo.
\end{example}

\begin{remark}[Interludio: el tamaño como biyección]\label{rem:b1:logic:interlude}
Un tema callado de este capítulo merece un foco: las biyecciones son la
noción matemática de «mismo tamaño». Para los \omterm{def:b1:logic:sets}{conjuntos} finitos se
convierte en el cálculo combinatorio del~\cref{ch:b1:counting}, donde
toda fórmula es en secreto una \omterm{def:b1:logic:inj}{biyección}; para los \omterm{def:b1:logic:sets}{conjuntos} infinitos,
en el problema del fin de semana que cierra el capítulo, donde $\N$,
$\Q$ y $\R$ resultan tener tamaños genuinamente distintos. El mismo
diccionario reaparece dos veces más en este volumen, en formas
refinadas: las sucesiones (\cref{ch:b1:seq}) no son otra cosa que
aplicaciones $\N \to \R$, de modo que los \omterm{def:b1:logic:statement}{enunciados} sobre
sucesiones son \omterm{def:b1:logic:statement}{enunciados} sobre un \omterm{def:b1:logic:sets}{conjunto} de aplicaciones; y el
álgebra lineal medirá los espacios vectoriales no con biyecciones, sino
con biyecciones \emph{lineales}, cuya existencia está gobernada por un
único número, la dimensión (\cref{ch:b1:findim}). Siempre que aparece
una nueva «igualdad» --- equipotencia, isomorfismo de grupos
(\cref{ch:b1:structures}), isomorfismo lineal --- se repite el patrón
del~\cref{thm:b1:logic:inverse}: la igualdad es una \omterm{def:b1:logic:map}{aplicación}
invertible que respeta la estructura.
\end{remark}

\begin{remark}[Dónde se usa este capítulo]\label{rem:b1:logic:whereused}
En todas partes --- pero algunos lugares merecen señalarse. La gimnasia
de tres cuantificadores del~\cref{ex:b1:logic:limit} es el pan de cada
día de los \cref{ch:b1:seq,ch:b1:continuity}: toda demostración de un
límite es una partida jugada contra un $\varepsilon$ arbitrario. Las
\omterm{def:b1:logic:equiv}{clases de equivalencia} reaparecen como las clases de congruencia de
$\Z/n\Z$ en el~\cref{ch:b1:structures}, donde la \omterm{thm:b1:logic:partition}{partición}
del~\cref{thm:b1:logic:partition} adquiere una estructura algebraica
propia. Las relaciones de \omterm{def:b1:logic:order}{orden}, las cotas superiores y los extremos
superiores se convierten en el corazón axiomático de $\R$ en
el~\cref{ch:b1:reals}. Las inyecciones, sobreyecciones y biyecciones
vuelven como las aplicaciones lineales del~\cref{ch:b1:linmaps}, donde
la \omterm{def:b1:logic:inj}{inyectividad} se puede comprobar sobre un solo vector (el núcleo); y
el problema del fin de semana convierte la noción desnuda de \omterm{def:b1:logic:inj}{biyección}
en una teoría de los \emph{tamaños de los \omterm{def:b1:logic:sets}{conjuntos} infinitos}, cuyas
conclusiones (numerabilidad de $\Q$, no numerabilidad de $\R$)
reaparecen en los \cref{ch:b1:reals,ch:b1:topology}.
\end{remark}

\section{Ejercicios}

\begin{exercise}[$\star$]\label{exo:b1:logic:1}
Escríbase la negación de cada \omterm{def:b1:logic:statement}{enunciado} sin emplear la palabra «no»:
\begin{enumerate}
  \item $\forall x \in \R,\ \exists y \in \R,\ x + y > 0$;
  \item $\exists x \in \R,\ \forall y \in \R,\ xy = 0$;
  \item $\forall \varepsilon > 0,\ \exists \delta > 0,\ \forall x \in
        \R,\ \abs{x} \leq \delta \implies \abs{f(x)} \leq \varepsilon$
        (para una \omterm{def:b1:logic:map}{aplicación} fija $f \colon \R \to \R$).
\end{enumerate}
Decídase después si los \omterm{def:b1:logic:statement}{enunciados} (1) y (2) son verdaderos.
\end{exercise}

\begin{exercise}[$\star$]\label{exo:b1:logic:2}
Sean $P, Q$ dos \omterm{def:b1:logic:statement}{enunciados}. Demuéstrese con tablas de verdad que
$\lnot(P \implies Q) \iff P \land (\lnot Q)$, y dedúzcase la negación
de: «si una función es derivable, entonces es continua».
\end{exercise}

\begin{exercise}[$\star$]\label{exo:b1:logic:3}
Demuéstrese por contraposición: para $x \in \R$, si $x^3 + x \geq 2$,
entonces $x \geq 1$. Demuéstrese después por reducción al absurdo que no
existe el menor número real estrictamente positivo.
\end{exercise}

\begin{exercise}[$\star$]\label{exo:b1:logic:4}
Demuéstrese por inducción que, para todo $n \in \N$:
\begin{enumerate}
  \item $\sum_{k=0}^{n} 2^k = 2^{n+1} - 1$;
  \item $4^n + 5$ es divisible por $3$.
\end{enumerate}
\end{exercise}

\begin{exercise}[$\star$]\label{exo:b1:logic:5}
Encuéntrese el fallo de la siguiente «demostración» de que todos los
lápices tienen el mismo color. \emph{Sea $P(n)$: «en todo \omterm{def:b1:logic:sets}{conjunto} de
$n$ lápices, todos los lápices tienen el mismo color». $P(1)$ es
evidente. Supongamos $P(n)$ y tomemos $n+1$ lápices; quitando el último,
los $n$ primeros comparten color; quitando el primero, los $n$ últimos
comparten color; luego los $n+1$ comparten color.}
\end{exercise}

\begin{exercise}[$\star$]\label{exo:b1:logic:6}
Sean $A, B, C$ subconjuntos de $E$. Demuéstrese:
\begin{enumerate}
  \item $A \setminus B = A \cap \overline{B}$;
  \item $(A \cup B) \setminus C = (A \setminus C) \cup (B \setminus
        C)$;
  \item $A \subseteq B \iff A \cup B = B \iff A \cap B = A$.
\end{enumerate}
\end{exercise}

\begin{exercise}[$\star\star$]\label{exo:b1:logic:7}
Decídase, con demostración, si cada \omterm{def:b1:logic:map}{aplicación} es \omterm{def:b1:logic:inj}{inyectiva},
\omterm{def:b1:logic:inj}{sobreyectiva} o \omterm{def:b1:logic:inj}{biyectiva}:
\begin{enumerate}
  \item $f \colon \N \to \N$, $n \mapsto n + 1$;
  \item $g \colon \Z \to \Z$, $n \mapsto n + 1$;
  \item $h \colon \R \setminus \{1\} \to \R$, $x \mapsto
        \frac{x+1}{x-1}$.
\end{enumerate}
Para $h$, ajústese el codominio para hacerla \omterm{def:b1:logic:inj}{biyectiva} y calcúlese la
inversa.
\end{exercise}

\begin{exercise}[$\star\star$]\label{exo:b1:logic:8}
Sea $f \colon E \to F$, y sean $A, A' \subseteq E$ y $B, B' \subseteq
F$.
\begin{enumerate}
  \item Demuéstrese que $f^{-1}(B \cap B') = f^{-1}(B) \cap f^{-1}(B')$
        y que $f(A \cup A') = f(A) \cup f(A')$.
  \item Demuéstrese que $f(A \cap A') \subseteq f(A) \cap f(A')$ y dese
        un ejemplo en el que la inclusión sea estricta.
  \item Demuéstrese que $f$ es \omterm{def:b1:logic:inj}{inyectiva} si y solo si $f(A \cap A') =
        f(A) \cap f(A')$ para todos $A, A'$.
\end{enumerate}
\end{exercise}

\begin{exercise}[$\star\star$]\label{exo:b1:logic:9}
Sean $f \colon E \to F$ y $g \colon F \to E$ tales que $g \circ f =
\mathrm{id}_E$. Demuéstrese que $f$ es \omterm{def:b1:logic:inj}{inyectiva} y $g$ es
\omterm{def:b1:logic:inj}{sobreyectiva}. Dese un ejemplo en el que ni $f$ ni $g$ sea \omterm{def:b1:logic:inj}{biyectiva}.
\end{exercise}

\begin{exercise}[$\star\star$]\label{exo:b1:logic:10}
En $\R$ se define $x \mathbin{\mathcal{R}} y \iff x^2 - y^2 = x - y$.
Demuéstrese que $\mathcal{R}$ es una \omterm{def:b1:logic:equiv}{relación de equivalencia} y
descríbase la \omterm{def:b1:logic:equiv}{clase de equivalencia} de cada real $x$. ¿Qué clases
tienen exactamente un elemento?
\end{exercise}

\begin{exercise}[$\star\star\star$]\label{exo:b1:logic:11}
(Cantor) Sea $E$ un \omterm{def:b1:logic:sets}{conjunto}. Demuéstrese que no hay ninguna
\omterm{def:b1:logic:inj}{sobreyección} de $E$ sobre $\mathcal{P}(E)$. \emph{Indicación: dada
$f \colon E \to \mathcal{P}(E)$, considérese $D = \{x \in E : x \notin f(x)\}$.}
\end{exercise}

\begin{exercise}[$\star\star\star$]\label{exo:b1:logic:12}
Sea $f \colon E \to F$ una \omterm{def:b1:logic:map}{aplicación}. Defínase
$\Phi \colon \mathcal{P}(F) \to \mathcal{P}(E)$ por $\Phi(B) = f^{-1}(B)$.
\begin{enumerate}
  \item Demuéstrese que $f$ es \omterm{def:b1:logic:inj}{sobreyectiva} si y solo si $\Phi$ es
        \omterm{def:b1:logic:inj}{inyectiva}.
  \item Demuéstrese que $f$ es \omterm{def:b1:logic:inj}{inyectiva} si y solo si $\Phi$ es
        \omterm{def:b1:logic:inj}{sobreyectiva}.
\end{enumerate}
\end{exercise}

\section{Problema: comparar infinitos}

\begin{problem}\label{pb:b1:logic:1}
¿Cuándo tienen dos \omterm{def:b1:logic:sets}{conjuntos} «el mismo número de elementos»? La
respuesta de Cantor --- cuando existe una \omterm{def:b1:logic:inj}{biyección} entre ellos ---
resulta ser utilizable incluso para \omterm{def:b1:logic:sets}{conjuntos} infinitos, y parte el
infinito en tamaños genuinamente distintos. Este problema construye toda
la caja de herramientas a partir de las definiciones desnudas del
capítulo: el teorema de Cantor--Schröder--Bernstein (dos inyecciones
fabrican una \omterm{def:b1:logic:inj}{biyección}), la numerabilidad de $\Q$, la no numerabilidad
de $\R$ por el argumento diagonal y la asombrosa conclusión de Cantor en
1874: \emph{existen \omterm{pb:b1:logic:1}{números trascendentes}, y en cantidad masiva}, sin
exhibir ni uno solo.
\index{conjunto numerable}\index{teorema de Cantor--Schröder--Bernstein}
En todo el problema, para \omterm{def:b1:logic:sets}{conjuntos} $E$ y $F$, se escribe $E \preceq F$
cuando existe una \omterm{def:b1:logic:inj}{inyección} de $E$ en $F$, y $E \approx F$ («$E$ y $F$
son \emph{equipotentes}») cuando existe una \omterm{def:b1:logic:inj}{biyección} de $E$ sobre $F$.

\medskip
\noindent\textbf{Parte I --- El vocabulario de la comparación.}
\begin{enumerate}
  \item Pruébese que $\approx$ se comporta como una \omterm{def:b1:logic:equiv}{relación de
        equivalencia}: $E \approx E$; si $E \approx F$, entonces
        $F \approx E$; si $E \approx F$ y $F \approx G$, entonces
        $E \approx G$. (Cítense con precisión
        el~\cref{thm:b1:logic:inverse} y la~\cref{prop:b1:logic:comp}.)
  \item Pruébese que $\preceq$ es transitiva y que una \omterm{def:b1:logic:inj}{inyección}
        $f \colon E \to F$ induce siempre $E \approx f(E)$.
  \item Sea $E \neq \emptyset$. Pruébese que $E \preceq F$ si y solo si
        existe una \omterm{def:b1:logic:inj}{sobreyección} de $F$ sobre $E$.
  \item Compruébese que $n \mapsto n + 1$ es una \omterm{def:b1:logic:inj}{biyección} de $\N$ sobre
        $\N^* = \N \setminus \{0\}$, y que
        \[
        \sigma(n) = \frac n2 \ \ (n \text{ par}), \qquad
        \sigma(n) = -\frac{n+1}2 \ \ (n \text{ impar})
        \]
        es una \omterm{def:b1:logic:inj}{biyección} de $\N$ sobre $\Z$. Así pues, quitar un punto,
        o duplicar hacia los negativos, no cambia el tamaño de $\N$.
\end{enumerate}

\medskip
\noindent\textbf{Parte II --- El teorema de
Cantor--Schröder--Bernstein.} Sean $f \colon E \to F$ y
$g \colon F \to E$ dos inyecciones. Defínanse
\[
C_0 = E \setminus g(F), \qquad C_{n+1} = g\bigl(f(C_n)\bigr)
\ \ (n \in \N), \qquad C = \bigcup_{n \in \N} C_n,
\]
y sea $h \colon E \to F$ la \omterm{def:b1:logic:map}{aplicación} que envía $x \in C$ a $f(x)$, y
$x \notin C$ al único $y \in F$ tal que $g(y) = x$.
\begin{enumerate}[resume]
  \item Compruébese que $h$ está bien definida: si $x \notin C$,
        entonces $x \in g(F)$, y el elemento $y$ con $g(y) = x$ es único.
  \item Pruébese que $g\bigl(f(C)\bigr) = \bigcup_{n \geq 1} C_n
        \subseteq C$. (Las imágenes directas conmutan con las uniones:
        \cref{exo:b1:logic:8}.)
  \item Pruébese que $h$ es \omterm{def:b1:logic:inj}{inyectiva}. (Tres casos; en el caso mixto
        $x \in C$, $x' \notin C$, véase que $h(x) = h(x')$ obligaría a
        que $x' \in g(f(C)) \subseteq C$.)
  \item Pruébese que $h$ es \omterm{def:b1:logic:inj}{sobreyectiva}: dado $y \in F$, distínganse
        los casos $g(y) \notin C$ y $g(y) \in C_n$ para algún $n
        \geq 1$ (¿por qué es imposible $g(y) \in C_0$?), y exhíbase una
        \omterm{def:b1:logic:map}{imagen recíproca} de $y$ en cada caso.
  \item Conclúyase el \emph{teorema de
        Cantor--Schröder--Bernstein}: si $E \preceq F$ y $F \preceq E$,
        entonces $E \approx F$. Coméntese en una frase qué hace que este
        \omterm{def:b1:logic:statement}{enunciado} no sea trivial.
  \item Dos aplicaciones. (a) Pruébese que $\intcc01 \approx \intoo01$.
        (b) Pruébese que $\varphi(p, q) = 2^p(2q + 1) - 1$ define una
        \omterm{def:b1:logic:inj}{biyección} de $\N \times \N$ sobre $\N$ --- la \omterm{def:b1:logic:inj}{inyectividad} por un
        argumento de paridad, la \omterm{def:b1:logic:inj}{sobreyectividad} por inducción fuerte
        (\cref{thm:b1:logic:induction}). Así pues, $\N \times \N \approx
        \N$: el plano de los puntos enteros no es mayor que la recta.
\end{enumerate}

\medskip
\noindent\textbf{Parte III --- \omterm{pb:b1:logic:1}{Conjuntos numerables}.} Un \omterm{def:b1:logic:sets}{conjunto} $E$ se
llama \emph{a lo sumo numerable} cuando $E \preceq \N$, y
\emph{numerable} cuando $E \approx \N$.
\begin{enumerate}[resume]
  \item Pruébese que todo subconjunto infinito $A \subseteq \N$ es
        numerable. (Defínase $\varphi(n)$ por recurrencia como el mínimo
        de $A \setminus \{\varphi(0), \dots, \varphi(n-1)\}$; véase que
        $\varphi$ es estrictamente creciente, que cumple
        $\varphi(n) \geq n$ y que alcanza todos los elementos de $A$.)
  \item Dedúzcase que un \omterm{def:b1:logic:sets}{conjunto} es a lo sumo numerable si y solo si
        es finito o numerable, y obsérvese que la pregunta 9 da el
        atajo: si $E \preceq \N$ y $\N \preceq E$, entonces $E$ es
        numerable.
  \item Pruébese que si $E$ y $F$ son a lo sumo numerables, también lo
        es $E \times F$. Dedúzcase que $\Z \times \N^*$ es numerable.
  \item Pruébese que $\Q$ es numerable. (Inyéctese $\Q$ en $\Z \times
        \N^*$ escribiendo cada racional como fracción irreducible de
        denominador positivo --- la unicidad de esa representación se
        demuestra en el~\cref{ch:b1:arith}; aplíquese después la pregunta 12.)
  \item Pruébese que una unión numerable de \omterm{pb:b1:logic:1}{conjuntos numerables} es a
        lo sumo numerable: si cada $E_n$ ($n \in \N$) es a lo sumo
        numerable, también lo es $\bigcup_{n \in \N} E_n$. (Envíese $x$
        al par $(n, f_n(x))$, donde $n$ es el índice \emph{mínimo} con
        $x \in E_n$.)
  \item Pruébese que el \omterm{def:b1:logic:sets}{conjunto} de los subconjuntos \emph{finitos} de
        $\N$ es numerable. (Envíese un subconjunto finito $F$ a
        $\sum_{i \in F} 2^i$; demuéstrese la \omterm{def:b1:logic:inj}{inyectividad} comparando el
        mayor elemento en el que difieren dos \omterm{def:b1:logic:sets}{conjuntos} finitos, con
        ayuda de $\sum_{k=0}^{m-1} 2^k = 2^m - 1$
        del~\cref{exo:b1:logic:4}.)
\end{enumerate}

\medskip
\noindent\textbf{Parte IV --- Diagonalización.} Denótese por
$\{0,1\}^{\N}$ el \omterm{def:b1:logic:sets}{conjunto} de todas las aplicaciones
$u \colon \N \to \{0, 1\}$, es decir, el \omterm{def:b1:logic:sets}{conjunto} de las sucesiones
binarias.
\begin{enumerate}[resume]
  \item Constrúyase una \omterm{def:b1:logic:inj}{biyección} entre $\mathcal{P}(\N)$ y
        $\{0,1\}^{\N}$ (funciones indicadoras).
  \item (El argumento diagonal) Sea $\Phi \colon \N \to
        \{0,1\}^{\N}$ una \omterm{def:b1:logic:map}{aplicación} cualquiera. Considérese la
        sucesión $d$ definida por $d(n) = 1 - \Phi(n)(n)$. Pruébese que
        $d$ no está en la imagen de $\Phi$, y conclúyase que
        $\{0,1\}^{\N}$ \emph{no} es a lo sumo numerable. Explíquese en
        una frase por qué, a través de la pregunta 17, esto es
        exactamente el teorema de Cantor (\cref{exo:b1:logic:11}) para
        $E = \N$.
  \item Admítase --- como es familiar desde la secundaria y se establece
        con rigor en el~\cref{ch:b1:reals} --- que todo $x \in
        \intco01$ tiene un único desarrollo decimal \emph{propio} $x =
        0.d_1 d_2 d_3\dots$ (uno que no termine en una sucesión infinita
        de cifras $9$). Dada una sucesión cualquiera $(x_n)_{n \geq 1}$ de
        elementos de $\intco01$, constrúyase $x \in \intco01$ con
        $x \neq x_n$ para todo $n$: tómese como $n$-ésima cifra un $5$ si
        la $n$-ésima cifra de $x_n$ es distinta de $5$, y un $6$ en caso
        contrario. Justifíquese con cuidado que $x$ es propio y que
        evita todos los $x_n$, y conclúyase que $\intco01$ no es a lo
        sumo numerable.
  \item Dedúzcase que $\R$ no es numerable, y que el \omterm{def:b1:logic:sets}{conjunto}
        $\R \setminus \Q$ de los números irracionales tampoco lo es. ¿En
        qué sentido preciso «casi todos» los números reales son
        irracionales?
\end{enumerate}

\medskip
\noindent\textbf{Parte V --- El teorema de Cantor de 1874: existen
\omterm{pb:b1:logic:1}{números trascendentes}.} Un número real $x$ es \emph{algebraico}
\index{número algebraico} cuando $P(x) = 0$ para algún polinomio no nulo
$P$ con coeficientes enteros, y \emph{trascendente}
\index{número trascendente} en caso contrario. Admítase en esta parte
--- se demuestra en el~\cref{ch:b1:poly} --- que un polinomio no nulo de
grado $n$ tiene a lo sumo $n$ raíces reales.
\begin{enumerate}[resume]
  \item Pruébese que todo número racional es algebraico, y hállense
        polinomios explícitos con coeficientes enteros que anulen a
        $\sqrt 2$ y a $\sqrt 2 + \sqrt 3$.
  \item Para $n \in \N$ fijo, pruébese que el \omterm{def:b1:logic:sets}{conjunto} de los polinomios
        de grado a lo sumo $n$ con coeficientes enteros es numerable.
        (Inyéctese en $\Z^{n+1}$ y aplíquese inducción sobre $n$ con la
        pregunta 13.)
  \item Dedúzcase que el \omterm{def:b1:logic:sets}{conjunto} de \emph{todos} los polinomios con
        coeficientes enteros es numerable.
  \item Demuéstrese el \emph{teorema de Cantor sobre los
        \omterm{pb:b1:logic:1}{números algebraicos}}: el \omterm{def:b1:logic:sets}{conjunto} $\mathcal{A}$ de los números
        reales algebraicos es numerable.
  \item Conclúyase: existen números reales trascendentes, y el
        \omterm{def:b1:logic:sets}{conjunto} de los \omterm{pb:b1:logic:1}{números trascendentes} no es numerable. Hágase
        después balance de todo el problema en unas pocas frases: la
        cadena $\N \approx \Z \approx \Q \approx \mathcal{A}$, el salto
        estricto hasta $\R \approx$ (esencialmente) $\mathcal{P}(\N)$,
        dónde fue decisiva cada herramienta
        (Cantor--Schröder--Bernstein, uniones numerables, la diagonal)
        --- y la fuerza filosófica de demostrar que existen incontables
        \omterm{pb:b1:logic:1}{números trascendentes} sin nombrar ni uno. (Demostrar que un
        número \emph{concreto} como $\pi$ es trascendente es mucho más
        difícil y queda fuera de este volumen.)
\end{enumerate}
\end{problem}
